% Options for packages loaded elsewhere
\PassOptionsToPackage{unicode}{hyperref}
\PassOptionsToPackage{hyphens}{url}
\documentclass[
]{article}
\usepackage{xcolor}
\usepackage[margin=1in]{geometry}
\usepackage{amsmath,amssymb}
\setcounter{secnumdepth}{-\maxdimen} % remove section numbering
\usepackage{iftex}
\ifPDFTeX
  \usepackage[T1]{fontenc}
  \usepackage[utf8]{inputenc}
  \usepackage{textcomp} % provide euro and other symbols
\else % if luatex or xetex
  \usepackage{unicode-math} % this also loads fontspec
  \defaultfontfeatures{Scale=MatchLowercase}
  \defaultfontfeatures[\rmfamily]{Ligatures=TeX,Scale=1}
\fi
\usepackage{lmodern}
\ifPDFTeX\else
  % xetex/luatex font selection
\fi
% Use upquote if available, for straight quotes in verbatim environments
\IfFileExists{upquote.sty}{\usepackage{upquote}}{}
\IfFileExists{microtype.sty}{% use microtype if available
  \usepackage[]{microtype}
  \UseMicrotypeSet[protrusion]{basicmath} % disable protrusion for tt fonts
}{}
\makeatletter
\@ifundefined{KOMAClassName}{% if non-KOMA class
  \IfFileExists{parskip.sty}{%
    \usepackage{parskip}
  }{% else
    \setlength{\parindent}{0pt}
    \setlength{\parskip}{6pt plus 2pt minus 1pt}}
}{% if KOMA class
  \KOMAoptions{parskip=half}}
\makeatother
\usepackage{color}
\usepackage{fancyvrb}
\newcommand{\VerbBar}{|}
\newcommand{\VERB}{\Verb[commandchars=\\\{\}]}
\DefineVerbatimEnvironment{Highlighting}{Verbatim}{commandchars=\\\{\}}
% Add ',fontsize=\small' for more characters per line
\usepackage{framed}
\definecolor{shadecolor}{RGB}{248,248,248}
\newenvironment{Shaded}{\begin{snugshade}}{\end{snugshade}}
\newcommand{\AlertTok}[1]{\textcolor[rgb]{0.94,0.16,0.16}{#1}}
\newcommand{\AnnotationTok}[1]{\textcolor[rgb]{0.56,0.35,0.01}{\textbf{\textit{#1}}}}
\newcommand{\AttributeTok}[1]{\textcolor[rgb]{0.13,0.29,0.53}{#1}}
\newcommand{\BaseNTok}[1]{\textcolor[rgb]{0.00,0.00,0.81}{#1}}
\newcommand{\BuiltInTok}[1]{#1}
\newcommand{\CharTok}[1]{\textcolor[rgb]{0.31,0.60,0.02}{#1}}
\newcommand{\CommentTok}[1]{\textcolor[rgb]{0.56,0.35,0.01}{\textit{#1}}}
\newcommand{\CommentVarTok}[1]{\textcolor[rgb]{0.56,0.35,0.01}{\textbf{\textit{#1}}}}
\newcommand{\ConstantTok}[1]{\textcolor[rgb]{0.56,0.35,0.01}{#1}}
\newcommand{\ControlFlowTok}[1]{\textcolor[rgb]{0.13,0.29,0.53}{\textbf{#1}}}
\newcommand{\DataTypeTok}[1]{\textcolor[rgb]{0.13,0.29,0.53}{#1}}
\newcommand{\DecValTok}[1]{\textcolor[rgb]{0.00,0.00,0.81}{#1}}
\newcommand{\DocumentationTok}[1]{\textcolor[rgb]{0.56,0.35,0.01}{\textbf{\textit{#1}}}}
\newcommand{\ErrorTok}[1]{\textcolor[rgb]{0.64,0.00,0.00}{\textbf{#1}}}
\newcommand{\ExtensionTok}[1]{#1}
\newcommand{\FloatTok}[1]{\textcolor[rgb]{0.00,0.00,0.81}{#1}}
\newcommand{\FunctionTok}[1]{\textcolor[rgb]{0.13,0.29,0.53}{\textbf{#1}}}
\newcommand{\ImportTok}[1]{#1}
\newcommand{\InformationTok}[1]{\textcolor[rgb]{0.56,0.35,0.01}{\textbf{\textit{#1}}}}
\newcommand{\KeywordTok}[1]{\textcolor[rgb]{0.13,0.29,0.53}{\textbf{#1}}}
\newcommand{\NormalTok}[1]{#1}
\newcommand{\OperatorTok}[1]{\textcolor[rgb]{0.81,0.36,0.00}{\textbf{#1}}}
\newcommand{\OtherTok}[1]{\textcolor[rgb]{0.56,0.35,0.01}{#1}}
\newcommand{\PreprocessorTok}[1]{\textcolor[rgb]{0.56,0.35,0.01}{\textit{#1}}}
\newcommand{\RegionMarkerTok}[1]{#1}
\newcommand{\SpecialCharTok}[1]{\textcolor[rgb]{0.81,0.36,0.00}{\textbf{#1}}}
\newcommand{\SpecialStringTok}[1]{\textcolor[rgb]{0.31,0.60,0.02}{#1}}
\newcommand{\StringTok}[1]{\textcolor[rgb]{0.31,0.60,0.02}{#1}}
\newcommand{\VariableTok}[1]{\textcolor[rgb]{0.00,0.00,0.00}{#1}}
\newcommand{\VerbatimStringTok}[1]{\textcolor[rgb]{0.31,0.60,0.02}{#1}}
\newcommand{\WarningTok}[1]{\textcolor[rgb]{0.56,0.35,0.01}{\textbf{\textit{#1}}}}
\usepackage{graphicx}
\makeatletter
\newsavebox\pandoc@box
\newcommand*\pandocbounded[1]{% scales image to fit in text height/width
  \sbox\pandoc@box{#1}%
  \Gscale@div\@tempa{\textheight}{\dimexpr\ht\pandoc@box+\dp\pandoc@box\relax}%
  \Gscale@div\@tempb{\linewidth}{\wd\pandoc@box}%
  \ifdim\@tempb\p@<\@tempa\p@\let\@tempa\@tempb\fi% select the smaller of both
  \ifdim\@tempa\p@<\p@\scalebox{\@tempa}{\usebox\pandoc@box}%
  \else\usebox{\pandoc@box}%
  \fi%
}
% Set default figure placement to htbp
\def\fps@figure{htbp}
\makeatother
\setlength{\emergencystretch}{3em} % prevent overfull lines
\providecommand{\tightlist}{%
  \setlength{\itemsep}{0pt}\setlength{\parskip}{0pt}}
\usepackage{bookmark}
\IfFileExists{xurl.sty}{\usepackage{xurl}}{} % add URL line breaks if available
\urlstyle{same}
\hypersetup{
  pdftitle={Bayesian Hypothesis Testing},
  hidelinks,
  pdfcreator={LaTeX via pandoc}}

\title{Bayesian Hypothesis Testing}
\author{}
\date{\vspace{-2.5em}2026-04-17}

\begin{document}
\maketitle

\subsection{Introduction}\label{introduction}

Bayesian decision-making about effects splits into two cultural
traditions that complement rather than replace each other. The
Bayes-factor school computes a ratio of marginal likelihoods and reads
strength of evidence from Jeffreys' scale. The Region of Practical
Equivalence (ROPE) school, popularised by Kruschke, defines a range of
effects that are practically null and asks how much of the posterior
lies inside it. Both produce decisions about hypotheses, but they answer
slightly different questions: Bayes factors compare models, while ROPE
asks whether the data have ruled out a substantively negligible effect.

\subsection{Prerequisites}\label{prerequisites}

A working understanding of posterior distributions, credible intervals,
marginal likelihoods, and the difference between statistical
significance and practical importance.

\subsection{Theory}\label{theory}

\textbf{ROPE.} Specify, before looking at the data, a parameter range
\([a, b]\) inside which an effect is considered practically zero --- for
example \(\beta \in (-0.1, 0.1)\) on a standardised scale. After
fitting, compute the proportion of the 95 \% HPD (or the full posterior)
inside that range. Conventional rules: \(> 95 \%\) inside means accept
the null for practical purposes, \(0 \%\) inside means reject it, and
intermediate values are inconclusive.

\textbf{Bayes factor.} Ratio of marginal likelihoods under two models;
for a point null nested in a continuous alternative, the Savage-Dickey
density ratio gives the same answer as the ratio of posterior to prior
densities at the null. Both quantities depend on the prior under the
alternative, so disclosing the prior is part of the report.

\subsection{Assumptions}\label{assumptions}

ROPE assumes a domain-relevant equivalence margin set in advance;
chasing margins after seeing the posterior is a Bayesian variant of
\(p\)-hacking. Bayes factors assume proper priors and a likelihood that
admits stable marginal-likelihood computation.

\subsection{R Implementation}\label{r-implementation}

\begin{Shaded}
\begin{Highlighting}[]
\FunctionTok{library}\NormalTok{(bayestestR); }\FunctionTok{library}\NormalTok{(brms)}

\FunctionTok{set.seed}\NormalTok{(}\DecValTok{2026}\NormalTok{)}
\NormalTok{d }\OtherTok{\textless{}{-}} \FunctionTok{data.frame}\NormalTok{(}\AttributeTok{y =} \FunctionTok{rnorm}\NormalTok{(}\DecValTok{200}\NormalTok{), }\AttributeTok{x =} \FunctionTok{rnorm}\NormalTok{(}\DecValTok{200}\NormalTok{))}
\NormalTok{fit }\OtherTok{\textless{}{-}} \FunctionTok{brm}\NormalTok{(y }\SpecialCharTok{\textasciitilde{}}\NormalTok{ x, }\AttributeTok{data =}\NormalTok{ d, }\AttributeTok{chains =} \DecValTok{4}\NormalTok{, }\AttributeTok{iter =} \DecValTok{2000}\NormalTok{, }\AttributeTok{refresh =} \DecValTok{0}\NormalTok{)}

\FunctionTok{rope}\NormalTok{(fit, }\AttributeTok{range =} \FunctionTok{c}\NormalTok{(}\SpecialCharTok{{-}}\FloatTok{0.1}\NormalTok{, }\FloatTok{0.1}\NormalTok{))}
\FunctionTok{bayesfactor\_parameters}\NormalTok{(fit)}
\FunctionTok{describe\_posterior}\NormalTok{(fit)}
\end{Highlighting}
\end{Shaded}

\subsection{Output \& Results}\label{output-results}

\texttt{rope()} returns the percentage of the HPD (or full posterior,
depending on the chosen \texttt{ci}) inside the equivalence region for
each parameter. \texttt{bayesfactor\_parameters()} returns the
Savage-Dickey Bayes factor against the point null.
\texttt{describe\_posterior()} is a one-stop summary that combines the
two with point estimates and credible intervals.

\subsection{Interpretation}\label{interpretation}

A reporting sentence: ``The posterior for \(\beta_x\) had 4 \% of its 95
\% HPD inside the practical-equivalence region \((-0.1, 0.1)\) and a
Savage-Dickey Bayes factor of \(\mathrm{BF}_{10} = 7.2\) --- substantial
evidence for a non-null effect on Jeffreys' scale, with most of the
posterior lying outside the practical-null region.'' When the two
methods disagree (ROPE inconclusive but Bayes factor decisive, or vice
versa), report both and explain which question your study was designed
to answer.

\subsection{Practical Tips}\label{practical-tips}

\begin{itemize}
\tightlist
\item
  Pre-register the ROPE based on what is practically relevant in your
  field; for standardised effects, \(\pm 0.1\) is a common default but
  is arbitrary in absolute terms.
\item
  For ROPE on raw scales, work out the smallest effect that would change
  clinical or business behaviour and use that as the half-width.
\item
  Bayes factors are sensitive to the prior under the alternative
  (Lindley--Bartlett paradox); always run a sensitivity analysis with at
  least two prior scales.
\item
  Use \texttt{bayestestR::describe\_posterior()} to deliver a single
  block of summaries that includes credible intervals, ROPE, Bayes
  factors, and probability of direction; consistent reporting reduces
  ambiguity.
\item
  ROPE is easier to communicate to non-statistical audiences (``most of
  the posterior is outside the negligible-effect range'') than Bayes
  factors are; choose the framing that fits your reader.
\end{itemize}

\subsection{R Packages Used}\label{r-packages-used}

\texttt{bayestestR} for \texttt{rope()},
\texttt{bayesfactor\_parameters()}, and \texttt{describe\_posterior()};
\texttt{brms} or \texttt{rstanarm} for the underlying Bayesian fit;
\texttt{bridgesampling} for marginal-likelihood-based Bayes factors when
the Savage-Dickey shortcut does not apply.

\subsection{For Reviewers}\label{for-reviewers}

\textbf{What to look for in a paper using this method.}

\begin{itemize}
\tightlist
\item
  \textbf{Common misapplications.}
\item
  \textbf{Diagnostics that should be reported but often aren't.}
\item
  \textbf{Red flags in tables and figures.}
\item
  \textbf{What to verify.}
\item
  \textbf{What an adequate Methods paragraph must contain.}
\end{itemize}

\subsection{See also --- labs in this
chapter}\label{see-also-labs-in-this-chapter}

\begin{itemize}
\tightlist
\item
  \href{labs/lab-bayesian-thinking-control-vs-decision-error.Rmd}{Bayesian
  thinking; α control vs decision error}
\item
  \href{labs/lab-brms-stan-loo-hierarchical-models.Rmd}{brms / Stan,
  LOO, hierarchical models}
\end{itemize}

\end{document}
